\documentclass[10pt,a4paper]{amsart}
\usepackage[margin=19mm]{geometry}
\usepackage{amsmath,amssymb,mathtools}
\usepackage[scr=rsfs,cal=euler]{mathalfa}
\usepackage{xfrac}
\usepackage[unicode,hidelinks,bookmarks=false]{hyperref}
\newcommand{\ie}{i.e.\ }
\newcommand{\cf}{cf.\ }
\newcommand{\resp}{respectively\ }
\newcommand{\Subsection}{\subsection}
\providecommand{\nobreakdash}{\nobreak\hbox{-}}
\setlength{\emergencystretch}{2em}
\multlinegap0pt
\usepackage{xfrac}
\newtheorem{lem}{Lemma}
\newcommand{\Ad}{\operatorname{Ad}} %Adjoint group action
\newcommand{\Span}{\operatorname{Span}} %Linear span
\newcommand{\g}{\mathfrak{g}} %Lie algebra
\newcommand{\tsp}[1]{{}^t \! #1} %French transpose
\title{Stratification of the Helffer-Nourrigat Cone}
\author{Cl\'ement Cren}
\date{Source: arXiv:2512.11434v2 (CC BY 4.0)}
\begin{document}
\maketitle

\noindent\emph{Source and licence.} Stratification of the Helffer-Nourrigat Cone. Version arXiv:2512.11434v2 (CC BY 4.0), distributed by arXiv under the Creative Commons Attribution 4.0 International licence, http://creativecommons.org/licenses/by/4.0/, as recorded in the arXiv licence metadata for this exact version and shown on the abstract page https://arxiv.org/abs/2512.11434v2. Full licence notice: \texttt{licenses/arxiv-2512.11434-CC-BY-4.0.txt}.\par\smallskip

\noindent\textit{Editorial scope.} Counted unit: the article's lem lem (source0.tex lines 547--549). The article's complete original proof of it is delivered with it (source0.tex lines 551--563, one \texttt{proof} environment of 13 lines). No mathematics is written, repaired or re-derived: the statement and the proof are the article's own lines, in the article's own notation, and no retained line is substituted or re-flowed.  No further source-local context is needed: every cross-reference of every retained line resolves inside the two delivered spans.  Nothing was omitted from any retained span, so the package carries no omission record.  No retained line contains a citation, so the wrapper carries no bibliography and no bibliography entry.  No retained line contains a figure environment, an image-inclusion command or drawing code, so no image is delivered and the package declares no asset dependency.  Editorial formatting only, in addition: an AMS \texttt{amsart} preamble that restates the article's own declarations and macros used by the retained lines, namely the environments \texttt{lem} on separate counters, together with the standard packages those lines need.  The retained environments print in this document as Lemma 1 in order of appearance, whereas the article prints the same items with the numbering of its own full document.  The complete native source of the article as distributed in the arXiv e-print for this exact version is bundled unchanged as \texttt{sources/190870/source0.tex}.
\begin{lem}
 Let $\g,\mathfrak{h}$ be graded Lie algebras, $\varphi\colon \g\twoheadrightarrow\mathfrak{h}$ a~surjective homomorphism that preserves the grading. Let $Y_1, \dots, Y_N$ be a~Jordan--H\"older basis of $\g$ compatible with the grading. Then there exists a~Jordan--H\"older basis of $\mathfrak{h}$ such that $\widehat{\varphi} \colon \widehat{\mathfrak{h}} \hookrightarrow \widehat{\g}$ maps elements belonging to a~same Pedersen stratum in $\widehat{\mathfrak{h}}$ to the same Pedersen stratum in $\widehat{\g}$. Moreover, the induced map between the Pedersen strata of $\widehat{\mathfrak{h}}$ and $\widehat{\g}$ is injective.
\end{lem}

\begin{proof}
 Define $Z_j = \varphi(Y_j)$, $1\leq j \leq N$. These vectors span the Lie algebra $\mathfrak{h}$. We extract a~basis in the following way. Define $i_1 = \min\{i\mid Z_i \neq 0\}$,
 and then inductively $i_{k+1} = \min\{i \mid Z_i \notin \Span(Z_{i_1},\dots, Z_{i_k})\}$,
 until we get a~basis for some $i_n$.
 The vectors $Z_{i_1}, \dots, Z_{i_n}$ form a~Jordan--H\"older basis of $\mathfrak{h}$, compatible with the grading.

 We now want to relate jump indices for the two bases when we apply $\tsp\varphi$. To do that, we define
$u \colon  \{1,\dots,N\}  \to \{1,\dots,n\}$, $ i \mapsto \max\{k\mid i_k \leq i\}$.
 The map $u$ is a~right inverse to $\sigma \colon k \mapsto i_k$. We have $\varphi(\g_i) = \mathfrak{h}_{u(i)}$ for all $1\leq i \leq N$. Recall also that if $\xi \in \mathfrak{h}^*$ and $g \in \g$ then \smash{$\tsp\varphi\bigl(\Ad^*_{\varphi(g)}(\xi)\bigr) = \Ad^*_g\bigl(\tsp\varphi(\xi)\bigr)$},
 i.e., $\tsp\varphi$ induces a~homeomorphism between the $\mathfrak{h}$-coadjoint orbit of $\xi$ and the $\g$-coadjoint orbit of~$\tsp\varphi(\xi)$.
 Consequently, for all $\xi \in \mathfrak{h}^*$ and $1\leq i \leq N$ we have $J^i\bigl(\tsp(\varphi(\xi)\bigr) = \sigma\bigl(J^{u(i)}(\xi)\bigr)$.
 We therefore map Pedersen strata to Pedersen strata, injectively between the strata.
\end{proof}

\end{document}
